\documentclass[12pt]{article}
\usepackage{amsmath}    
\usepackage{amssymb}   
\usepackage{setspace} 
\usepackage{geometry}   
\usepackage{booktabs}   
\usepackage{graphicx}  
\usepackage{hyperref}  
\usepackage{natbib}    
\usepackage{appendix}
\usepackage{ragged2e}
\usepackage{placeins}
\usepackage[style=apa,backend=biber]{biblatex}
\addbibresource{references.bib}
\usepackage{enumitem} 
\geometry{margin=1in}
\setstretch{1.2}
\usepackage{xcolor}
\title{DMP D2 (End of Semester Report)}
\author{Leyao Wang  pkv4su@virginia.edu}
\date{\today}
\begin{document}

\maketitle

\iffalse
\begin{centering}
\section*{Abstract (temp)}
\textcolor{red}{TBD}
\end{centering}
\fi

\section{Motivation}

Generative artificial intelligence (AI), particularly in the form of large language model (LLM)-integrated assistants, has undergone rapid development and diffusion since late 2022. Unlike previous waves of automation and technologies that primarily affected routine or low-skill work, the latest iteration of AI has demonstrated the capability to influence a wide spectrum of occupations, including those requiring advanced cognitive or analytical skills. A growing body of research focused on measuring associated productivity gains and efficiency improvements of incorporating AI, but empirical studies on the labor market's response to this technological shock remain limited. 

Theoretically, the impact of AI on the labor market presents an inherent ambiguity. On one hand, AI-powered tools may augment productivity by assisting with daily tasks, thereby remaining complementary to human labor. On the other hand, as LLMs become increasingly capable of performing certain tasks, they might substitute for labor in those segments. This dual mechanism reflects the distinction emphasized by \cite{acemoglu2019automation}: technological adoption can both displace and reinstate labor, depending on which tasks are automated, and which remain complementary to human expertise. 

These dynamics motivate the central question of this project: How has the development of generative AI reshaped labor-market outcomes across different experience levels and occupations? In particular, this study seeks to understand whether the diffusion of LLM-based tools has altered the relative economic advantages of professional experience knowledge-intensive work. 

\section{Research Questions}

This project focuses on two primary dimensions of analysis to empirically evaluate the impact of LLM-based tools: wages and occupational mobility. The analysis focuses on testing whether the rapid adoption of AI has reshaped these key labor-market outcomes in ways that differ between younger, mid-career, and senior professionals.

\subsection{Wage Premium of Professional Experience}

The first and primary dimension of analysis focuses on wage, which offers the most direct and observable measure of AI's labor market impact. Specifically, the project will test whether the introduction of AI assistants has altered the wage premium associated with experience within occupations highly exposed to these technologies. 

\begin{itemize}
    \item \textbf{Hypothesis A (Substitution)}: The wage premium for mid-career and senior workers has compressed relative to that of younger entrants in AI-exposed occupations. This outcome would be consistent with the partial substitution of accumulated experience, as AI systems reduce the comparative advantage of human expertise in performing certain specialized tasks. 
    \item \textbf{Hypothesis B (Complementarity):} The experience premium has widened, with wage gains accruing disproportionately to senior professionals.  One possible explanation is that senior professionals have successfully integrated AI-powered tools into their workflow and harnessed relatively more productivity gains comparing to junior workers. Another possible explanation is that the relative importance of senior developers’ work, for example, in oversight, system design, and mentorship, has become more significant in the post-AI period, enhancing their relative returns.
\end{itemize}

\subsection{Occupational Mobility}

The second dimension of analysis examines occupational mobility, assessing whether AI exposure has changed movement patterns between occupations or industries. Rather than focusing solely on complete labor force exit, this analysis will emphasize transitions between roles that differ in their degree of AI complementarity. This analysis aims to capture how workers adapt to technological change, whether voluntarily or involuntarily, via occupation change. For instance, workers facing potential wage premium compression may exhibit greater incentives to transition into new roles that are less exposed to AI's influence. 




\iffalse
\textcolor{red}{I consider dropping motivation and moving section RQ up here}
\noindent
Generative artificial intelligence, particularly large language model (LLM) – integrated assistants, has developed and diffused rapidly since late 2022. Unlike previous automation technologies that primarily affected routine or low-skill work, the latest wave of AI has influenced a wide range of occupations, including those requiring certain cognitive or analytical skills. LLM-based tools have been incorporated into professional workflows across different industries, and a growing body of research contributes to measuring productivity gains and efficiency improvements associated with AI. However, empirical studies on AI’s labor market impact remains limited.

\medskip
\noindent
Theoretically, AI’s impact on labor market is ambiguous. On one hand, AI powered tools may augment productivity by assisting with daily work. On the other hand, as LLMs become increasingly capable of delegating work, they may be substitute for certain tasks, reducing the demand for human labor in those segments. This dual mechanism reflects the distinction emphasized by \cite{acemoglu2019automation}: technological adoption can both displace and reinstate labor, depending on which tasks are automated, and which remain complementary to human expertise. 

\medskip
\noindent
These dynamics motivate the central question of my project: how has the rise of generative AI reshaped the labor-market outcomes across experience levels and occupations? In particular, I seek to understand whether the diffusion of LLM-based tools has altered the relative advantages of experience and specialization in knowledge-intensive work. 
\begin{itemize}
    \item \textbf{If the relationship between AI and professional experience is complementary:} Then we would expect wage gains to accrue disproportionately to senior professionals. One possible explanation is that senior professionals have successfully integrated AI-powered tools into their workflow and harnessed relatively more productivity gains comparing to junior workers. Another possible explanation is that the relative importance of senior developers’ work, for example, in oversight, system design, and mentorship, has become more significant in the post-AI period, enhancing their relative returns. 
    \item \textbf{If the relationship between AI and professional experience is substitutive:} Then mid-career workers and senior professionals may face compressed wage premia. They might also have greater incentives to transition into new roles that are less exposed to AI's influence. 
\end{itemize}

\textcolor{red}{*Seems that I could formalize this into a Model section. This will look cool...}
\section{Model}
\textcolor{red}{Planned. To be done before winter break}
%If AI complements to experience, enhancing the productivity of workers who can effectively integrate new tools into their work, then we should expect wage gains to accrue disproportionately to senior professionals. Conversely, if AI substitutes for accumulated expertise by automating the tasks, mid-career workers may face compressed wage premiums or greater incentives to transition into new roles.
%Understanding this distinction is crucial for interpreting how emerging technologies reshape high-skill labor markets and for forecasting longer-term changes in skill demand, occupational mobility, and career trajectories. 

\medskip
\noindent

%By systematically linking occupational exposure to generative AI with changes in wages and job movement patterns, this project aims to provide empirical evidence on how exposure to generative AI is transforming
\fi


\section{Partial Literature Review} 

\iffalse
\subsection{Disruptive Technological Shock}

\textcolor{red}{*TBD*}
\fi

\iffalse
\subsection{Task-Based Approaches}\textcolor{red}{Update needed}
\fi
\begin{itemize}[leftmargin=*]
    \item \textbf{(\cite{felten2018method})}  
    Introduces a quantitative approach linking progress in distinct AI subfields (from the EFF AI Progress metrics) to human abilities defined in O*NET.  
    Maps AI capabilities $\rightarrow$ O*NET abilities $\rightarrow$ occupations. 

    \item \textbf{(\cite{felten2021aioe})}  
    Extends the 2018 framework to a forward-looking exposure index using Mechanical Turk surveys to map ten AI applications to 52 O*NET abilities.  
    Aggregates ability-level exposure (weighted by importance and prevalence) to occupation, industry, and geographic levels.  
    Produces standardized exposure scores such as \textit{AI Occupational Exposure (AIOE)}, \textit{AI Industry Exposure (AIIE)}, and \textit{AI Geography Exposure (AIGE)}, measuring potential technical exposure rather than realized adoption.

    \item \textbf{(\cite{webb2020impact})}  
    Constructs a task-technology overlap model matching verbs and nouns from patent text to O*NET task descriptions.  
    Identifies high-skill cognitive tasks (judgment, analysis) as most AI-exposed, contrasting with robotics’ focus on routine manual work.  
    Provides a scalable, data-driven mapping of AI technologies to tasks using textual similarity.

    \item \textbf{(\cite{eloundou2024gpts})}  
    Uses GPT-4 and human annotators to classify O*NET tasks into categories based on whether LLMs could reduce task completion time by at least 50\%. Aggregates task labels to occupation-level exposure scores and produces composite indices.  
    Provides a large-scale, LLM-specific exposure dataset reflecting technical feasibility rather than observed adoption.
    \item \textbf{(\cite{acemoglu2019automation})}  
    Develops a theoretical model explaining how automation both displaces existing tasks and creates new ones.  
    Highlights the ambiguous effects of technology on labor demand depending on whether new tasks are complements or substitutes for human labor.  
\end{itemize}


\subsection{Empirical Evidence on AI's Labor Market Impact}
\iffalse
\textcolor{red}{Update needed}
\fi
\begin{itemize}[leftmargin=*]

    \item \textbf{(\cite{hosseini2025seniority})}  
Examines whether generative AI constitutes a seniority-biased technological change, disproportionately affecting junior relative to senior workers.  
Combines LinkedIn résumé data (Revelio Labs) with nearly 200 million job postings. Measures firm-level AI adoption by identifying postings for ``GenAI integrator'' roles responsible for implementing or integrating AI tools. Employs a difference-in-differences framework comparing junior versus senior employment in adopting versus non-adopting firms before and after 2023.  
Finds that after ChatGPT’s release, junior employment in adopting firms declined by roughly 9–10\%, while senior employment remained stable, with effects concentrated in high-exposure occupations.  Provides early empirical evidence that generative AI adoption shifts firm workforce composition toward senior labor.
\item \textbf{(\cite{brynjolfsson2025canaries})}  
Uses high-frequency administrative payroll data from ADP, covering millions of U.S. workers, to document six stylized facts about the early labor market effects of generative AI.  
Identifies that employment among early-career workers (ages 22–25) in highly AI-exposed occupations has declined by roughly 13\% since late 2022, even after controlling for firm-level shocks.  
Employment for older or less AI-exposed workers remained stable or continued to grow.  
Shows that employment declines are concentrated in occupations where AI automates rather than augments work, based on task-level measures from the Anthropic Economic Index and GPT-4 exposure (Eloundou et al., 2024).  
Finds limited wage effects, suggesting adjustments occurred primarily through employment rather than compensation.  
\end{itemize}

\iffalse
\section{Research Questions}
\textcolor{red}{UPDATE to background + RQ + transition to Lit Review}

This project will empirically examine how the development and diffusion of generative AI—particularly large language model (LLM)–based tools—have affected the careers of workers across different occupations and experience levels. The central question is whether the rapid adoption of these technologies has reshaped key labor-market outcomes, such as wages and occupational mobility, in ways that differ between younger and mid-career or senior professionals.

The first and primary dimension of analysis focuses on wages, which offer the most direct and observable measure of AI’s impact. Specifically, the project will test whether the introduction of AI-assisted tools has altered the wage premium associated with experience within occupations. One hypothesis is that the wage premium for mid-career workers has compressed relative to that of younger entrants in occupations highly exposed to AI technologies, consistent with partial substitution of accumulated experience as AI systems reduce the advantage of expertise. An alternative hypothesis is that the experience premium has widened if senior workers have effectively integrated AI tools into their workflows, thereby amplifying productivity. Another possible explanation is that the growing importance of oversight, system design, leadership, and mentorship roles in the post-AI period has enhanced the relative returns to experience.


The second dimension of analysis examines occupational mobility, assessing whether AI exposure has changed movement between occupations or industries. Rather than focusing solely on complete labor-force exit, the analysis will emphasize transitions between roles that differ in their degree of AI complementarity—such as shifts from technical to managerial positions. This dimension of analysis aims to capture how workers, voluntarily or involuntarily, adapt to technological change via occupation change.
\fi

\iffalse
\section{Model}
\textcolor{red}{Planned. Update needed}
\fi
\section{Data}
\iffalse
\textcolor{red}{Update needed}
\fi
The analysis will draw on multiple datasets, grouped into three main categories corresponding to the project’s objectives:\textbf{(1) Measures of AI Exposure, (2) Wage Outcomes, and (3) Occupational Mobility.}

\subsection{Measures of AI Exposure}
%\noindent \textbf{Measures of AI Exposure.} 

To quantify variation in exposure to large language models (LLMs) and related generative AI technologies across occupations, I plan to use the occupational-level exposure measures developed by \cite{eloundou2024gpts} in the paper \textit{``GPTs are GPTs: An Early Look at the Labor Market Impact of Large Language Models''}. The dataset, publicly available at \href{https://github.com/openai/GPTs-are-GPTs/tree/main/data}{https://github.com/openai/GPTs-are-GPTs}, provides exposure scores for U.S. occupations based on the \textit{Standard Occupational Classification (SOC)} system.  

These exposure scores are derived by linking task-level data from O*NET to assessments of whether each task could be performed or augmented by LLMs. The authors define three main indices: \(E_1\) (direct exposure), \(E_2\) (indirect exposure through task dependencies), and a combined measure \(E_1 + 0.5E_2\), which is used as a balanced indicator of overall LLM exposure. The \(E_1 + 0.5E_2\) composite will serve as the baseline measure of occupational exposure to generative AI in this project.  

This dataset is particularly suitable for my analysis because it captures exposure specifically to generative AI tools that emerged after 2022, aligning closely with the time frame of widespread adoption of AI-integrated assistants in professional work.

In future robustness checks, I may also incorporate the earlier \textit{AI Occupational Exposure (AIOE)} and \textit{AI Industry Exposure (AIIE)} indices (\cite{felten2021aioe}) to assess whether results are consistent when using pre-generative-AI exposure measures. Datasets on AIOE and AIIE are publicly available at \href{https://github.com/AIOE-Data/AIOE}{https://github.com/AIOE-Data/AIOE}.
\subsection{Wage Outcomes}
  
The main analysis of wage outcomes will rely on the \textit{American Community Survey (ACS)} and, as a supplementary source, the \textit{Current Population Survey (CPS)}.  
\begin{itemize}[leftmargin=*]
    \item The ACS and CPS-ASEC provide large annual samples containing information on wage income, occupation, age, education, and other demographic characteristics. It is well-suited for estimating wage differentials and returns to experience across occupations. However, due to the government shutdown, the release of ACS 2024 dataset is delayed. The preliminary results presented in this report is gathered from CPS-ASEC dataset - which is already available for 2024 and 2025.
    \item The CPS offers monthly observations on employment status, occupation, wages, and demographics. Although smaller in sample size, it could allow for higher-frequency tracking and serves as a useful robustness check.  
\end{itemize}

Both ACS and CPS use standardized occupational coding systems compatible with SOC codes, and they can be accessed directly through IPMUS. 
\subsection{Occupational Mobility}
%\noindent \textbf{Occupational Mobility.}
For studying mobility patterns, the CPS will serve as the primary data source. Its short panel structure allows for limited individual tracking, enabling the identification of short-term occupation switches or exits from a given occupation. While the CPS cannot capture long-term career outcomes, it provides the best available evidence on job mobility within and across occupations in a given period of time.  

As a backup option or for robustness checks, I may incorporate the \textit{Quarterly Census of Employment and Wages} dataset. Although (to the best of my knowledge) this dataset does not include occupational codes, it provides quarterly employment statistics aggregated at the industry level, which could complement the CPS-based analysis.

In summary, these datasets collectively allow for an integrated analysis linking LLM exposure (\cite{eloundou2024gpts}) with wage and mobility outcomes across the U.S. labor market. 


\section{Empirical Design \& Preliminary Results}
\subsection{Scope and Sample}
The empirical analysis in this report uses \textbf{CPS-ASEC} data accessed from IPMUS portal. I choose to focus on the college-educated labor force, as this group is presumably most directly exposed to the diffusion of generative AI technologies in professional and knowledge-intensive occupations and allows for greater comparability in educational background. The goal of the design is to estimate how the returns to expertise in the field —reflected as the \textit{tenure premium} in wages—have evolved across occupations with different levels of LLM exposure.

\subsection{Overview of the Empirical Strategy}
%\noindent \textbf{Overview of the Empirical Strategy.}  

\textcolor{blue}{Note: As pointed out by Prof. Harrington and Prof. Cosar, my current estimating approach can be, and should be updated to a more rigorous (single-stage) regression form with interactive terms and year dummies. I have definitely put this on my list, but due to time constraints, this update is \textbf{not} featured in this report.}


To analyze how AI diffusion has altered the wage structure, I propose a two-stage estimation framework. The first stage estimates, within each occupation and year, the relationship between wages and labor-market experience using a Mincer regression. The second stage then relates changes in the estimated experience premium across years to the exposure scores ($E_1+0.5E_2$) from \textbf{\cite{eloundou2024gpts}}'s dataset.

\medskip

\noindent \textbf{Stage 1: Estimating the return to experience.}  

The first stage uses the following Mincer equation to capture how wages vary with experience for worker $i$ in occupation $o$ during year (or quarter) $t$:
\[
\ln(w_{i,o,t}) = \beta_{0,o,t} + \beta_{1,o,t} Exp_{i} + \beta_2 Exp_{i}^2 + Z_{i} + \epsilon,
\]
\textit{where}
\begin{itemize}
    \item $\ln(w_{i,o,t})$ is the log wage of person $i$ in occupation $o$ in year $t$.
    \item $Exp_i$ is the person's years of experience in the field. 
    \item $Z_{i,t}$ is the set of vector of controls for personal background (geography, gender, highest degree earned, and ethnicity) 
\end{itemize}

The primary coefficient of interest is $\beta_{1}$, which represents the wage return to labor-market experience. I will also record $\beta_{2}$ to capture potential nonlinearities. 


\medskip
\noindent \textbf{Stage 2: Linking wage changes to LLM exposure.}

Next, I examine how the return to experience evolves with exposure to generative AI technologies. After selecting a baseline year $t_0$ (e.g. 2022), and the latest year $t_1$ of available data, I compute the change in the estimated experience premium for each occupation:
\[
\Delta \beta_{1,o} = \beta_{1,o,t_1} - \beta_{1,o,t_0}.
\]
This $\Delta \beta_{1,o}$ is then regressed on the LLM exposure score for occupation $o$:
\[
\Delta \beta_{1,o} = \gamma_0 + \gamma_1 Exposure_o + \epsilon,
\]
\textit{where}

\begin{itemize}
    \item $Exposure_o$ is the AI/LLM Exposure Score from \cite{eloundou2024gpts}
\end{itemize}

The coefficient $\gamma_1$ measures the association between exposure to LLMs and changes in the return to experience across occupations. A negative $\gamma_1$ would suggest that highly exposed occupations experienced a compression of the experience premium, consistent with LLMs substituting for accumulated expertise in a field; a positive $\gamma_1$ would indicate complementarity, where experience becomes more valuable after the difusion of AI-powered tools.

\medskip
\noindent \textbf{Advantages of the Two-Stage Approach.}  

This design has two main advantages.
\begin{itemize}
    \item It allows me to focus on how the return to experience change over time, without being affected by overall wage level differences across occupations.
    \item It provides a natural interpretation of AI's role in shaping the \textbf{career prospects} of professionals with different experience levels. Compared to a simple setup, my method better captures how recent diffusion of AI/LLM technology amplifies or erodes the relative value of accumulated experience in different occupations.
\end{itemize}

\textcolor{blue}{\textbf{One important limitation is that my approach CANNOT identify by which channel AI is impacting. Nor can my approach identify causality.} If AI is impacting work in different channels by both automating tasks (presumably more done by junior professionals) and creating new tasks (presumably more done by senior professionals) (\cite{acemoglu2019automation}), or they are leading senior roles to take up a larger portion in a team, then my observed result would be a composite of all these impacts.}

%First, it focuses directly on how the return to experience changes over time, without being affected by overall wage level differences across occupations. Second, it provides a natural interpretation of AI’s role in shaping the \textbf{career prospects} of workers with different experience levels. 
%Compared to a single-step regression of wages on AI exposure, this method better captures whether generative AI diffusion amplifies or erodes the relative value of accumulated experience in each occupation.

\medskip

\subsection{Preliminary Results}
As shown in \autoref{fig:b1-1}, college-educated workers are represented more in occupation groups with mid-to-high levels of LLM exposure. This pattern is consistent with my assumption that LLM technologies disproportionately affected occupations performed by more-educated workers. The general trend seems to suggest that there is \textbf{NO or VERY WEAK negative correlation} between LLM Exposure and returns to experience. 



\begin{figure}[h!]	
\caption{Stage 2 $\Delta\beta_1$ and LLM Exposure Score \label{fig:b1-1}}
\includegraphics[scale=.9]{graphs/12-03 on CPS-ASEC/ASEC_beta1_scat.pdf}
\includegraphics[scale=.9]{graphs/12-03 on CPS-ASEC/ASEC_beta1_scat_alt.pdf}
\justifying{\textbf{Notes:} Data source: CPS-ASEC. $t_1 = 2025$. $t_0 = 2022$. Occupations are bundled into 22 major groups as deifned by BLS.
Result is weighted by number of observations in each occupation group.}
\end{figure}

\FloatBarrier

Notably, the occupation group on the right end: \textbf{Computer and Mathematical Occupations}, seems to be an outlier that is moving in the opposite direction comparing to other mid-to-high levels LLM exposed occupation groups. 

To further investigate what happened to Computer Science and Mathematical Occupations, I plotted for the time trend for CS/Math and non-CS/Math occupations in \autoref{fig:b1-2}. Throughout the sample perido, the $\beta_1$ estimate for CS/Math occupations is consistently lower than that for other occupations. The gap widens between 2021 and 2024, before narrowing sharply in 2025. It is hard to say whether this turn is related to the disperse of LLM in any meaningful way, but this is something that is worth exploring in my future analysis.

\begin{figure}[h!]	
\caption{$\beta_1$ Estimates: CS/Math vs non-CS/Math Occupations \label{fig:b1-2}}
\includegraphics[scale=.9]{graphs/12-03 on CPS-ASEC/ASEC_beta1_trend_alt.pdf}
\justifying{\textbf{Notes:} Data source: CPS-ASEC. Result is weighted by number of observations in each group.}
\end{figure}

\FloatBarrier

For results related to $\beta_2$ estimates, please refer to \autoref{fig:b2-1} and \autoref{fig:b2-2} in the appendix. Because the estimated magnitude of $\beta_2$ is very small, its contribution to wage changes (for college-educated workers) is correspondingly limited relative to that of $\beta_1$

\clearpage

\subsection{Next Steps}
\noindent \textbf{Next Steps.}  
Several extensions are planned:
\begin{itemize}
    \item As suggested by Prof. Harrington and Prof. Cosar, update my empirical approach to a more rigorous regression form with interactive terms and year dummies. 
    \item Look into correlations between occupations and college major. I could be able to design an IV approach using worker's college major as the instrument for returns to experience. 
    \item As a robustness check, replicate the analysis using the earlier \textit{AI Occupational Exposure (AIOE)} and \textit{AI Industry Exposure (AIIE)} indices (\cite{felten2021aioe}) to test consistency across pre- and post-generative-AI exposure measures.
\end{itemize}


\section*{Limitations}

\begin{itemize}[leftmargin=*]
    \item I remain skeptical about the methodology used in \cite{eloundou2024gpts} to construct measures of AI exposure. The approach relies exclusively on human and GPT-4 annotators to classify O*NET tasks based on the likelihood that LLMs could perform or accelerate them, rather than on technical benchmarks or observed adoption data. As a result, the exposure measures capture potential feasibility rather than realized productivity or employment effects. \textcolor{blue}{To address this concern, I could use a more objective measure of AI adoption—the \textbf{Anthropic Economic Index}. Unlike Eloundou et al.’s subjective classification, this index leverages real-world Claude (a widely used LLM agent) conversation data to map actual AI usage onto O*NET tasks. This dataset is publicly available, though I have not yet examined its construction and coverage in detail.}

    \item The time horizon for observing AI’s labor-market impacts is short. With meaningful diffusion of generative AI beginning only in late 2022, the available data through 2025 provide at most three years of post-adoption observations.

    \item The ACS and CPS datasets are valuable due to their large sample sizes, standardized occupational codes, and consistent annual or monthly coverage. However, they are far from comprehensive, as they rely on survey samples rather than population-level administrative data. This sampling-based nature limits credibility of my results.
    \item My current approach would fail to account for most recent developments in AI tools 

\end{itemize}
\newpage
\begin{appendices}
\section{Appendix Figures}

\begin{figure}[h!]	
\caption{Stage 2 $\Delta\beta_2$ and LLM Exposure Score \label{fig:b2-1}}
\includegraphics[scale=.8]{graphs/12-03 on CPS-ASEC/ASEC_beta2_scat.pdf}
\includegraphics[scale=.8]{graphs/12-03 on CPS-ASEC/ASEC_beta2_scat_alt.pdf}
\justifying{\textbf{Notes:} Data source: CPS-ASEC. $t_1 = 2025$. $t_0 = 2022$. Occupations are bundled into 22 major groups as deifned by BLS.
Result is weighted by number of observations in each occupation group.}
\end{figure}

\FloatBarrier


\begin{figure}[h!]	
\caption{$\beta_2$ Estimates: CS/Math vs non-CS/Math Occupations \label{fig:b2-2}}
\includegraphics[scale=.9]{graphs/12-03 on CPS-ASEC/ASEC_beta2_trend_alt.pdf}
\justifying{\textbf{Notes:} Data source: CPS-ASEC. Result is weighted by number of observations in each group.}
\end{figure}



\end{appendices}

\clearpage
\printbibliography



\end{document}